\chapter{Επιλεγμένα Τμήματα Κώδικα}
\label{ch:appendix-code}

Ο πλήρης πηγαίος κώδικας της προσομοίωσης βρίσκεται στο αποθετήριο του
έργου, στους φακέλους \texttt{PROD/} (MATLAB) και \texttt{Model/}
(Python). Παρακάτω παρατίθενται τα πιο αντιπροσωπευτικά τμήματα, όπως
αναφέρονται στα Κεφάλαια~\ref{ch:methodology} και~\ref{ch:implementation}.

\section{Πυρήνας επιλογής κόμβου (\texttt{simulateScenario.m})}
\label{sec:appendix-simulatescenario}

Το παρακάτω απόσπασμα υλοποιεί τον κανόνα επιλογής κόμβου
(εξίσωση~\eqref{eq:node-selection}): για κάθε χρήστη, υπολογίζεται το SNR
προς κάθε υποψήφιο σταθμό βάσης (με στοχαστικό LOS draw και shadow fading,
εξισώσεις~\eqref{eq:los-prob-uma}-\eqref{eq:snr-bs}), κρατείται ο
καλύτερος υποψήφιος BS, υπολογίζεται ανεξάρτητα το SNR του δορυφόρου, και
οι δύο αξιολογούνται \emph{ο καθένας ξεχωριστά} έναντι του ελάχιστου
χρησιμοποιήσιμου SNR -- ο χρήστης συνδέεται σε όποιον/όποιους το
ξεπερνούν, ταυτόχρονα και στους δύο αν το ξεπερνούν και οι δύο
(πολυσυνδεσιμότητα).

\begin{lstlisting}[language=Matlab, caption={Απόσπασμα από \texttt{PROD/simulateScenario.m}: επιλογή κόμβου ανά χρήστη (πολυσυνδεσιμότητα).}]
for u = 1:numUsers
    userBestSNR  = -Inf;
    userBestNode = "";

    %% ===== Terrestrial BS candidates =====
    for b = 1:numBs
        % ... geodetic2enu, groundDistance, d3d ...

        % LOS probability, TR 38.901 par. 7.4.2
        pLos  = losProbability38901(groundDistance, user_geo(u,3), ...
                                     simParameters.PathLoss.Scenario);
        isLos = rand() < pLos;

        [pathLoss, sigmaSF] = nrPathLoss(simParameters.PathLoss, ...
                              simParameters.CarrierFrequency, isLos, ...
                              txPosition, rxPosition);

        % Shadow fading: log-normal, tipiki apoklisi sigmaSF
        pathLoss = pathLoss + sigmaSF * randn();

        % EIRP = TxPower + AntennaGain (TR 38.901 Pinakas 7.3-1)
        snr_db = (simParameters.EIRP - 30) - pathLoss - noisePowerBS_dBW;

        if snr_db > userBestSNR
            userBestSNR  = snr_db;
            userBestNode = "BS" + string(b);
        end
    end
    bsWinnerNode = userBestNode;   % kalyteros BS, prin sygkrithei me ton doryforo

    %% ===== Satellite candidate =====
    [~, elevSat, slantRangeSat] = geodetic2aer( ...
        sat_geo(1), sat_geo(2), sat_geo(3), ...
        user_geo(u,1), user_geo(u,2), user_geo(u,3), wgs84);

    if elevSat >= satParameters.MinElevationDeg
        lambdaSat = physconst('LightSpeed') / satParameters.CarrierFrequency;
        satPathLoss = fspl(slantRangeSat, lambdaSat);

        % Atmosfairiki aposvesi aerion, TR 38.811 ex. (6.6-8)
        satPathLoss = satPathLoss + gasAttenuationSlantP676( ...
            satParameters.CarrierFrequency, elevSat);

        satSnrDb = (satParameters.EIRP - 30) - satPathLoss - noisePowerSAT_dBW;
    else
        satPathLoss = inf;
        satSnrDb = -Inf;
    end

    %% ===== Apofasi syndesimotitas (SS-SBS, tesseris katastaseis) =====
    bsUsable  = userBestSNR >= minUsableSnrDb;
    satUsable = satSnrDb    >= minUsableSnrDb;

    if bsUsable && satUsable
        bestNodeTypeVec(u) = "DualConnectivity";
        bestNodeVec(u)     = bsWinnerNode + "+SAT-1";
    elseif bsUsable
        bestNodeTypeVec(u) = "Terrestrial";
        bestNodeVec(u)     = bsWinnerNode;
    elseif satUsable
        bestNodeTypeVec(u) = "Satellite";
        bestNodeVec(u)     = "SAT-1";
    else
        bestNodeTypeVec(u) = "Outage";
        bestNodeVec(u)     = "None";
    end
end
\end{lstlisting}

\section{Μοντέλο κίνησης LEO (\texttt{temporalPassSimulation.m})}
\label{sec:appendix-temporal}

Το παρακάτω απόσπασμα υλοποιεί τις εξισώσεις~\eqref{eq:orbital-period}-\eqref{eq:ground-track}:
τον υπολογισμό της Κεπλεριανής τροχιακής περιόδου/ταχύτητας ίχνους εδάφους,
και τον χρονικό βρόχο που μετακινεί το υποδορυφορικό σημείο και καλεί εκ
νέου τη \texttt{simulateScenario} σε κάθε βήμα, περνώντας την κατάσταση
καναλιού (\texttt{channelState}) από βήμα σε βήμα ώστε το shadow fading
να είναι χωρικά συσχετισμένο κατά Gudmundson
(εξίσωση~\eqref{eq:correlated-sf}) αντί να επαναδειγματίζεται ανεξάρτητα.

\begin{lstlisting}[language=Matlab, caption={Απόσπασμα από \texttt{PROD/temporalPassSimulation.m}: ταχύτητα ίχνους εδάφους και χρονικός βρόχος.}]
muEarth = 3.986004418e14;  % m^3/s^2
Re      = 6371e3;          % m, mesi aktina Gis
a       = Re + satAltitude;
orbitalPeriodS  = 2*pi*sqrt(a^3/muEarth);      % Kepleriani periodos (s)
groundSpeedMps  = (2*pi/orbitalPeriodS) * Re;  % taxytita ixnous edafous (m/s)

channelState = []; % kamia proigoumeni katastasi prin to proto vima
while step < maxSteps
    offsetKm = startOffsetKm + groundSpeedMps * t / 1000;
    subLon = centerLon + offsetKm / (111.320*cosd(centerLat));
    sat_geo = [centerLat, subLon, satAltitude];

    rng(step + 1);

    % to channelState perna vima-vima, xorika sysxetismeno (Gudmundson 1991)
    [bestNodeVec, bestNodeTypeVec, ~, ~, bestSnrDbVec, capacityMbpsVec, ...
        ~, ~, energyPerBitUJVec, ~, ~, ~, ~, satElevationVec, ~, ~, ...
        channelState] = ...
        simulateScenario(bs_geo, user_geo, sat_geo, wgs84, ...
                          simParameters, satParameters, channelState);

    % ... katagrafi apotelesmaton ana xristi ...

    refElev = max(satElevationVec);
    if refElev >= satParameters.MinElevationDeg
        wasVisible = true;
    elseif wasVisible
        break; % i dielefsi elixe
    end

    t = t + dtSeconds;
    step = step + 1;
end
\end{lstlisting}

\section{Χαρακτηριστικά και προεπεξεργασία μοντέλου μηχανικής μάθησης (\texttt{train\_model.py})}
\label{sec:appendix-train-model}

Το παρακάτω απόσπασμα ορίζει τα χαρακτηριστικά εισόδου (μόνο per-candidate
διαγνωστικά, όχι τα χαρακτηριστικά του ήδη-επιλεγμένου νικητή κόμβου -- βλ.
Ενότητα~\ref{sec:meth-ml}) και τον διαχωρισμό εκπαίδευσης/ελέγχου ανά σενάριο.

\begin{lstlisting}[language=Python, caption={Απόσπασμα από \texttt{Model/train\_model.py}: χαρακτηριστικά και grouped train/test split.}]
FEATURE_COLUMNS_NUMERIC = [
    # BsLoad/SatLoad exclude: synepeia tou ServingType, oxi predictor.
    "NumBS", "NumUsers",
    "CandBS_SNR_dB", "CandBS_Distance_m", "CandBS_PathLoss_dB",
    "CandSat_SNR_dB", "CandSat_Elevation_deg", "CandSat_SlantRange_m",
    "CandSat_PathLoss_dB",
]
# ServingType: Terrestrial/Satellite/DualConnectivity (3-class, vlepe ml_common.py).
TARGET_COLUMN = "ServingType"

def group_train_test_split(df, test_size=0.25, seed=42):
    # Split ana ScenarioID, oxi ana grammi - apofygi diarrois geometrias senariou.
    splitter = GroupShuffleSplit(n_splits=1, test_size=test_size,
                                  random_state=seed)
    train_idx, test_idx = next(splitter.split(df, groups=df["ScenarioID"]))
    return df.iloc[train_idx], df.iloc[test_idx]
\end{lstlisting}
